\chapter*{Annexes}
\addcontentsline{toc}{chapter}{Annexes}

\sloppy % Évite les débordements de ligne

\section*{Annexe A : Scripts SQL complets}
\addcontentsline{toc}{section}{Annexe A : Scripts SQL complets}

\subsection*{A.1 Création du schéma et des dimensions}
\addcontentsline{toc}{subsection}{A.1 Création du schéma et des dimensions}

\begin{verbatim}
-- ============================================================
-- SCRIPT COMPLET DE CRÉATION DE LA BASE DE DONNÉES
-- Base : biat_budget
-- Schéma : budget_biat
-- ============================================================

-- Création du schéma
CREATE SCHEMA budget_biat;
SET search_path TO budget_biat;

-- Dimension Année
CREATE TABLE dim_annee (
    id SERIAL PRIMARY KEY,
    annee INTEGER NOT NULL UNIQUE,
    date_debut DATE,
    date_fin DATE,
    est_bissextile BOOLEAN DEFAULT FALSE
);

-- Dimension Mois
CREATE TABLE dim_mois (
    id SERIAL PRIMARY KEY,
    mois_numero INTEGER NOT NULL UNIQUE,
    mois_nom VARCHAR(20) NOT NULL,
    mois_nom_court VARCHAR(3),
    trimestre INTEGER,
    semestre INTEGER,
    jours_mois INTEGER
);

-- Dimension Structure
CREATE TABLE dim_structure (
    id SERIAL PRIMARY KEY,
    code_structure VARCHAR(50) NOT NULL UNIQUE,
    nom_structure VARCHAR(200) NOT NULL,
    niveau VARCHAR(50),
    parent_id INTEGER,
    actif BOOLEAN DEFAULT TRUE
);

-- Dimension Employé
CREATE TABLE dim_employe (
    id SERIAL PRIMARY KEY,
    matricule VARCHAR(20) NOT NULL UNIQUE,
    nom VARCHAR(100),
    prenom VARCHAR(100),
    sexe CHAR(1),
    date_naissance DATE,
    date_embauche DATE,
    date_sortie DATE,
    categorie VARCHAR(50),
    statut VARCHAR(50),
    code_structure VARCHAR(50),
    actif BOOLEAN DEFAULT TRUE
);

-- Dimension Budget
CREATE TABLE dim_budget (
    id SERIAL PRIMARY KEY,
    annee INTEGER NOT NULL,
    version INTEGER DEFAULT 1,
    statut VARCHAR(30) DEFAULT 'PROVISOIRE',
    UNIQUE(annee, version)
);
\end{verbatim}

\subsection*{A.2 Création des tables de faits}
\addcontentsline{toc}{subsection}{A.2 Création des tables de faits}

\begin{verbatim}
-- Fait Charges Employé
CREATE TABLE fact_charges_employe (
    id SERIAL PRIMARY KEY,
    annee INTEGER NOT NULL,
    mois INTEGER NOT NULL,
    matricule VARCHAR(20) NOT NULL,
    code_structure VARCHAR(50),
    salaire_base DECIMAL(18,2) DEFAULT 0,
    anciennete DECIMAL(18,2) DEFAULT 0,
    treizieme_mois DECIMAL(18,2) DEFAULT 0,
    prime_fonction DECIMAL(18,2) DEFAULT 0,
    prime_individuelle_variable DECIMAL(18,2) DEFAULT 0,
    heures_supplementaires DECIMAL(18,2) DEFAULT 0,
    prime_resultat DECIMAL(18,2) DEFAULT 0,
    prime_exceptionnelle DECIMAL(18,2) DEFAULT 0,
    prime_mobilite DECIMAL(18,2) DEFAULT 0,
    tickets_restaurant DECIMAL(18,2) DEFAULT 0,
    mutuelle DECIMAL(18,2) DEFAULT 0,
    prevoyance DECIMAL(18,2) DEFAULT 0,
    transport DECIMAL(18,2) DEFAULT 0,
    creche_jardin_enfants DECIMAL(18,2) DEFAULT 0,
    oeuvres_sociales DECIMAL(18,2) DEFAULT 0,
    cotisations_sociales DECIMAL(18,2) DEFAULT 0,
    cotisation_retraite_complementaire DECIMAL(18,2) DEFAULT 0,
    frais_formation DECIMAL(18,2) DEFAULT 0,
    indemnite_depart DECIMAL(18,2) DEFAULT 0,
    indemnite_retraite DECIMAL(18,2) DEFAULT 0,
    solde_tout_compte DECIMAL(18,2) DEFAULT 0,
    frais_recrutement DECIMAL(18,2) DEFAULT 0,
    frais_integration DECIMAL(18,2) DEFAULT 0,
    total_brut DECIMAL(18,2) DEFAULT 0,
    total_charges_patronales DECIMAL(18,2) DEFAULT 0,
    total_global DECIMAL(18,2) DEFAULT 0,
    date_chargement TIMESTAMP DEFAULT NOW()
);

-- Fait Charges Structure
CREATE TABLE fact_charges_structure (
    id SERIAL PRIMARY KEY,
    annee INTEGER NOT NULL,
    mois INTEGER NOT NULL,
    code_structure VARCHAR(50),
    nb_employes INTEGER DEFAULT 0,
    total_charges DECIMAL(18,2) DEFAULT 0,
    date_chargement TIMESTAMP DEFAULT NOW()
);

-- Fait Budget Structure
CREATE TABLE fact_budget_structure (
    id SERIAL PRIMARY KEY,
    annee INTEGER NOT NULL,
    code_structure VARCHAR(50),
    id_budget INTEGER,
    budget_total DECIMAL(18,2) DEFAULT 0,
    date_chargement TIMESTAMP DEFAULT NOW()
);
\end{verbatim}

\subsection*{A.3 Création des tables de prédiction}
\addcontentsline{toc}{subsection}{A.3 Création des tables de prédiction}

\begin{verbatim}
-- Prédictions ML
CREATE TABLE pred_charges_mensuelles (
    id SERIAL PRIMARY KEY,
    annee INTEGER NOT NULL,
    mois INTEGER NOT NULL,
    code_structure VARCHAR(50),
    prediction_ml DECIMAL(18,2) NOT NULL,
    borne_inferieure DECIMAL(18,2),
    borne_superieure DECIMAL(18,2),
    niveau_confiance DECIMAL(5,2),
    modele_utilise VARCHAR(100),
    date_prediction TIMESTAMP DEFAULT NOW()
);

-- Écarts budgétaires
CREATE TABLE pred_ecarts_budgetaires (
    id SERIAL PRIMARY KEY,
    annee INTEGER NOT NULL,
    mois INTEGER NOT NULL,
    code_structure VARCHAR(50),
    consomme DECIMAL(18,2),
    budget_initial DECIMAL(18,2),
    prediction_ml DECIMAL(18,2),
    ecart_pourcentage DECIMAL(5,2),
    niveau_alerte VARCHAR(10),
    date_analyse TIMESTAMP DEFAULT NOW()
);

-- Risques de dépassement
CREATE TABLE pred_risque_depassement (
    id SERIAL PRIMARY KEY,
    annee INTEGER NOT NULL,
    mois INTEGER NOT NULL,
    code_structure VARCHAR(50),
    risque_depassement BOOLEAN,
    probabilite_risque DECIMAL(5,2),
    facteur_principal VARCHAR(200),
    date_analyse TIMESTAMP DEFAULT NOW()
);

-- Simulations
CREATE TABLE pred_simulations (
    id SERIAL PRIMARY KEY,
    annee INTEGER NOT NULL,
    mois INTEGER NOT NULL,
    code_structure VARCHAR(50),
    scenario VARCHAR(30),
    variation_effectif DECIMAL(5,2),
    variation_salaire DECIMAL(5,2),
    impact_charges DECIMAL(18,2),
    date_simulation TIMESTAMP DEFAULT NOW()
);
\end{verbatim}

\subsection*{A.4 Création des tables d'administration}
\addcontentsline{toc}{subsection}{A.4 Création des tables d'administration}

\begin{verbatim}
-- Table des administrateurs
CREATE TABLE admin_users (
    id SERIAL PRIMARY KEY,
    username VARCHAR(50) NOT NULL UNIQUE,
    password_hash VARCHAR(255) NOT NULL,
    email VARCHAR(100) NOT NULL UNIQUE,
    nom_complet VARCHAR(100),
    role VARCHAR(30) DEFAULT 'ADMIN',
    dernier_acces TIMESTAMP,
    date_creation TIMESTAMP DEFAULT NOW(),
    actif BOOLEAN DEFAULT TRUE
);

-- Table des sessions
CREATE TABLE sessions (
    id SERIAL PRIMARY KEY,
    token TEXT NOT NULL UNIQUE,
    admin_id INTEGER REFERENCES admin_users(id),
    date_creation TIMESTAMP DEFAULT NOW(),
    date_expiration TIMESTAMP,
    actif BOOLEAN DEFAULT TRUE
);

-- Insertion du compte admin par défaut
INSERT INTO admin_users (username, password_hash, email, nom_complet, role)
VALUES ('admin', 'admin123', 'admin@biat.com.tn', 'Administrateur BIAT', 'SUPER_ADMIN');
\end{verbatim}

\subsection*{A.5 Création des index pour optimisation}
\addcontentsline{toc}{subsection}{A.5 Création des index pour optimisation}

\begin{verbatim}
-- Index sur les clés étrangères
CREATE INDEX idx_charges_annee_mois ON fact_charges_employe(annee, mois);
CREATE INDEX idx_charges_matricule ON fact_charges_employe(matricule);
CREATE INDEX idx_charges_structure ON fact_charges_employe(code_structure);
CREATE INDEX idx_charges_structure_agg 
    ON fact_charges_structure(annee, mois, code_structure);
CREATE INDEX idx_budget_structure ON fact_budget_structure(annee, code_structure);
CREATE INDEX idx_pred_annee_mois ON pred_charges_mensuelles(annee, mois, code_structure);
CREATE INDEX idx_ecarts_alerte ON pred_ecarts_budgetaires(niveau_alerte);
CREATE INDEX idx_employe_matricule ON dim_employe(matricule);
CREATE INDEX idx_structure_code ON dim_structure(code_structure);
\end{verbatim}

\subsection*{A.6 Trigger pour calculs automatiques}
\addcontentsline{toc}{subsection}{A.6 Trigger pour calculs automatiques}

\begin{verbatim}
-- Trigger pour fact_charges_employe
CREATE OR REPLACE FUNCTION update_totaux_charges()
RETURNS TRIGGER AS $$
BEGIN
    -- Calcul du total brut
    NEW.total_brut := COALESCE(NEW.salaire_base,0) + COALESCE(NEW.anciennete,0) + 
                      COALESCE(NEW.treizieme_mois,0) + COALESCE(NEW.prime_fonction,0) +
                      COALESCE(NEW.prime_individuelle_variable,0) + 
                      COALESCE(NEW.heures_supplementaires,0) +
                      COALESCE(NEW.prime_resultat,0) + COALESCE(NEW.prime_exceptionnelle,0) +
                      COALESCE(NEW.prime_mobilite,0) + COALESCE(NEW.tickets_restaurant,0) +
                      COALESCE(NEW.mutuelle,0) + COALESCE(NEW.prevoyance,0) +
                      COALESCE(NEW.transport,0) + COALESCE(NEW.creche_jardin_enfants,0) +
                      COALESCE(NEW.oeuvres_sociales,0);
    
    -- Calcul du total global
    NEW.total_global := NEW.total_brut + COALESCE(NEW.cotisations_sociales,0) + 
                        COALESCE(NEW.cotisation_retraite_complementaire,0);
    RETURN NEW;
END;
$$ LANGUAGE plpgsql;

CREATE TRIGGER trg_update_totaux_charges
    BEFORE INSERT OR UPDATE ON fact_charges_employe
    FOR EACH ROW
    EXECUTE FUNCTION update_totaux_charges();
\end{verbatim}

\newpage

\section*{Annexe B : Code Python de l'API Flask (extraits clés)}
\addcontentsline{toc}{section}{Annexe B : Code Python de l'API Flask}

\subsection*{B.1 Configuration et connexion à la base de données}
\addcontentsline{toc}{subsection}{B.1 Configuration et connexion à la base de données}

\begin{verbatim}
# Configuration de l'application
app = Flask(__name__)
CORS(app, origins=["http://localhost:3000"])

app.config['JWT_SECRET_KEY'] = 'biat_budget_secret_key_2025'
app.config['JWT_ACCESS_TOKEN_EXPIRES'] = timedelta(hours=8)

DB_CONFIG = {
    'host': 'localhost',
    'port': 5432,
    'database': 'biat_budget',
    'user': 'postgres',
    'password': 'postgres'
}

def get_db_connection():
    return psycopg2.connect(**DB_CONFIG)
\end{verbatim}

\subsection*{B.2 Authentification JWT}
\addcontentsline{toc}{subsection}{B.2 Authentification JWT}

\begin{verbatim}
@app.route('/api/login', methods=['POST'])
def login():
    data = request.get_json()
    username = data.get('username')
    password = data.get('password')
    
    if username == 'admin' and password == 'admin123':
        access_token = create_access_token(identity={
            'id': 1,
            'username': 'admin',
            'role': 'ADMIN'
        })
        
        return jsonify({
            'success': True,
            'token': access_token,
            'user': {
                'id': 1,
                'username': 'admin',
                'nom_complet': 'Administrateur BIAT'
            }
        })
    
    return jsonify({'success': False, 'message': 'Identifiants invalides'}), 401

def token_required(f):
    @wraps(f)
    @jwt_required()
    def decorated(*args, **kwargs):
        current_user = get_jwt_identity()
        return f(current_user, *args, **kwargs)
    return decorated
\end{verbatim}

\subsection*{B.3 API de consultation des données}
\addcontentsline{toc}{subsection}{B.3 API de consultation des données}

\begin{verbatim}
@app.route('/api/data/<table_name>', methods=['GET'])
@token_required
def get_table_data(current_user, table_name):
    page = int(request.args.get('page', 1))
    per_page = int(request.args.get('per_page', 50))
    offset = (page - 1) * per_page
    
    conn = get_db_connection()
    cursor = conn.cursor(cursor_factory=RealDictCursor)
    
    cursor.execute(f"SELECT COUNT(*) as total FROM budget_biat.{table_name}")
    total = cursor.fetchone()['total']
    
    cursor.execute(f"""
        SELECT * FROM budget_biat.{table_name}
        ORDER BY id
        LIMIT %s OFFSET %s
    """, (per_page, offset))
    
    data = cursor.fetchall()
    
    return jsonify({
        'success': True,
        'data': data,
        'total': total,
        'page': page,
        'per_page': per_page,
        'total_pages': (total + per_page - 1) // per_page
    })
\end{verbatim}

\subsection*{B.4 API d'export des données}
\addcontentsline{toc}{subsection}{B.4 API d'export des données}

\begin{verbatim}
@app.route('/api/export/<table_name>', methods=['GET'])
@token_required
def export_table(current_user, table_name):
    format_type = request.args.get('format', 'csv')
    
    conn = get_db_connection()
    df = pd.read_sql(f"SELECT * FROM budget_biat.{table_name}", conn)
    conn.close()
    
    filename = f"{table_name}_{datetime.now().strftime('%Y%m%d_%H%M%S')}"
    
    if format_type == 'csv':
        output = io.StringIO()
        df.to_csv(output, index=False, encoding='utf-8-sig')
        output.seek(0)
        return send_file(
            io.BytesIO(output.getvalue().encode('utf-8-sig')),
            mimetype='text/csv',
            as_attachment=True,
            download_name=f"{filename}.csv"
        )
    
    return jsonify({'success': False, 'message': 'Format non supporté'}), 400
\end{verbatim}

\newpage

\section*{Annexe C : Code React des composants principaux}
\addcontentsline{toc}{section}{Annexe C : Code React des composants principaux}

\subsection*{C.1 Composant Login}
\addcontentsline{toc}{subsection}{C.1 Composant Login}

\begin{verbatim}
import React, { useState } from 'react';
import { useNavigate } from 'react-router-dom';

function Login() {
    const [username, setUsername] = useState('');
    const [password, setPassword] = useState('');
    const [error, setError] = useState('');
    const [loading, setLoading] = useState(false);
    const navigate = useNavigate();

    const handleSubmit = async (e) => {
        e.preventDefault();
        setLoading(true);
        setError('');

        try {
            const response = await fetch('http://localhost:5000/api/login', {
                method: 'POST',
                headers: { 'Content-Type': 'application/json' },
                body: JSON.stringify({ username, password })
            });

            const data = await response.json();

            if (data.success) {
                localStorage.setItem('token', data.token);
                localStorage.setItem('user', JSON.stringify(data.user));
                navigate('/dashboard');
            } else {
                setError(data.message);
            }
        } catch (err) {
            setError('Erreur de connexion au serveur');
        } finally {
            setLoading(false);
        }
    };

    return (
        <div style={styles.container}>
            <div style={styles.card}>
                <h1>🏦 BIAT Budget</h1>
                <form onSubmit={handleSubmit}>
                    <input type="text" placeholder="Nom d'utilisateur" 
                           value={username} 
                           onChange={(e) => setUsername(e.target.value)} />
                    <input type="password" placeholder="Mot de passe" 
                           value={password} 
                           onChange={(e) => setPassword(e.target.value)} />
                    {error && <div style={styles.error}>{error}</div>}
                    <button type="submit" disabled={loading}>
                        {loading ? 'Connexion...' : 'Se connecter'}
                    </button>
                </form>
            </div>
        </div>
    );
}

export default Login;
\end{verbatim}

\subsection*{C.2 Composant Dashboard (extrait)}
\addcontentsline{toc}{subsection}{C.2 Composant Dashboard (extrait)}

\begin{verbatim}
function Dashboard() {
    const [stats, setStats] = useState({});
    const navigate = useNavigate();

    useEffect(() => {
        fetchStats();
    }, []);

    const fetchStats = async () => {
        const token = localStorage.getItem('token');
        const response = await fetch('http://localhost:5000/api/dashboard/stats', {
            headers: { 'Authorization': `Bearer ${token}` }
        });
        const data = await response.json();
        if (data.success) setStats(data.stats);
    };

    return (
        <div>
            <div style={styles.header}>
                <h1>BIAT Budget</h1>
                <button onClick={() => navigate('/visualizations')}>
                    📊 Visualisations
                </button>
            </div>
            <div style={styles.statsGrid}>
                <div style={styles.statCard}>
                    <div>{stats.nb_employes || 0}</div>
                    <div>Employés</div>
                </div>
                {-- Autres cartes statistiques... --}
            </div>
        </div>
    );
}
\end{verbatim}

\newpage

\section*{Annexe D : Guide d'installation et déploiement}
\addcontentsline{toc}{section}{Annexe D : Guide d'installation et déploiement}

\subsection*{D.1 Prérequis système}
\addcontentsline{toc}{subsection}{D.1 Prérequis système}

\begin{itemize}
    \item \textbf{Système d'exploitation} : Windows 10/11, Linux (Ubuntu 20.04+), ou macOS
    \item \textbf{Mémoire RAM} : 8 Go minimum recommandé
    \item \textbf{Espace disque} : 5 Go disponibles
    \item \textbf{Processeur} : 2 cœurs minimum
\end{itemize}

\subsection*{D.2 Installation de PostgreSQL}
\addcontentsline{toc}{subsection}{D.2 Installation de PostgreSQL}

\begin{enumerate}
    \item Télécharger PostgreSQL 15 depuis \texttt{https://www.postgresql.org/download/}
    \item Exécuter l'installateur et suivre les instructions
    \item Mémoriser le mot de passe superutilisateur (postgres)
    \item Conserver le port par défaut (5432)
    \item Installer pgAdmin (inclus dans l'installateur)
\end{enumerate}

\subsection*{D.3 Installation de Python et des dépendances}
\addcontentsline{toc}{subsection}{D.3 Installation de Python et des dépendances}

\begin{verbatim}
# Télécharger Python 3.11 depuis python.org

# Créer un dossier pour le projet
mkdir biat_budget_app
cd biat_budget_app

# Créer un environnement virtuel
python -m venv venv

# Activer l'environnement virtuel
# Windows :
venv\Scripts\activate
# Linux/Mac :
source venv/bin/activate

# Installer les dépendances
pip install -r requirements.txt
\end{verbatim}

\subsection*{D.4 Configuration de la base de données}
\addcontentsline{toc}{subsection}{D.4 Configuration de la base de données}

\begin{verbatim}
# Se connecter à PostgreSQL
psql -U postgres

# Exécuter les scripts SQL dans l'ordre
\i scripts/01_create_schema.sql
\i scripts/02_create_tables.sql
\i scripts/03_create_indexes.sql
\i scripts/04_insert_data.sql

# Vérifier l'installation
SELECT COUNT(*) FROM budget_biat.dim_employe;
\end{verbatim}

\subsection*{D.5 Configuration de l'API Flask}
\addcontentsline{toc}{subsection}{D.5 Configuration de l'API Flask}

\begin{enumerate}
    \item Copier le fichier \texttt{app.py} dans le dossier \texttt{backend/}
    \item Modifier les paramètres de connexion dans \texttt{DB\_CONFIG}
    \item Exécuter l'API :
\end{enumerate}

\begin{verbatim}
cd backend
python app.py
\end{verbatim}

\subsection*{D.6 Configuration du frontend React}
\addcontentsline{toc}{subsection}{D.6 Configuration du frontend React}

\begin{verbatim}
# Installer Node.js 18 depuis nodejs.org

# Créer l'application React
npx create-react-app frontend
cd frontend

# Installer les dépendances
npm install react-router-dom axios

# Copier les composants dans src/components/
# Lancer l'application
npm start
\end{verbatim}

\subsection*{D.7 Vérification du bon fonctionnement}
\addcontentsline{toc}{subsection}{D.7 Vérification du bon fonctionnement}

\begin{enumerate}
    \item Vérifier que PostgreSQL tourne : \texttt{sudo systemctl status postgresql}
    \item Vérifier l'API : \texttt{curl http://localhost:5000/api/health}
    \item Accéder à l'application : \texttt{http://localhost:3000}
    \item Se connecter avec \texttt{admin / admin123}
\end{enumerate}

\newpage

\section*{Annexe E : Manuel utilisateur}
\addcontentsline{toc}{section}{Annexe E : Manuel utilisateur}

\subsection*{E.1 Connexion à l'application}
\addcontentsline{toc}{subsection}{E.1 Connexion à l'application}

\begin{enumerate}
    \item Ouvrir un navigateur web (Chrome, Firefox, Edge)
    \item Accéder à l'URL : \texttt{http://localhost:3000}
    \item Saisir les identifiants :
    \begin{itemize}
        \item Nom d'utilisateur : \texttt{admin}
        \item Mot de passe : \texttt{admin123}
    \end{itemize}
    \item Cliquer sur le bouton "Se connecter"
\end{enumerate}

% Espace pour figure - à remplacer par \includegraphics si nécessaire
% \begin{figure}[h!]
% \centering
% \caption{Page de connexion à l'application}
% \end{figure}

\subsection*{E.2 Navigation dans le Dashboard}
\addcontentsline{toc}{subsection}{E.2 Navigation dans le Dashboard}

Le dashboard présente quatre sections principales :

\begin{enumerate}
    \item \textbf{En-tête} : affiche le logo, l'heure et le bouton de déconnexion
    \item \textbf{Cartes statistiques (KPIs)} : 
    \begin{itemize}
        \item Nombre d'employés
        \item Nombre de structures
        \item Total des charges
        \item Meilleur modèle avec sa MAE
    \end{itemize}
    \item \textbf{Actions rapides} : accès direct aux pages principales
    \item \textbf{Données disponibles} : liste des tables consultables
\end{enumerate}

% Espace pour figure
% \begin{figure}[h!]
% \centering
% \caption{Dashboard principal avec les KPIs}
% \end{figure}

\subsection*{E.3 Consultation des données}
\addcontentsline{toc}{subsection}{E.3 Consultation des données}

\begin{enumerate}
    \item Cliquer sur une carte dans la section "Données disponibles"
    \item Utiliser la barre de recherche pour filtrer les données
    \item Naviguer entre les pages avec les boutons "Précédent" et "Suivant"
    \item Exporter les données via les boutons CSV, Excel ou JSON
\end{enumerate}

% Espace pour figure
% \begin{figure}[h!]
% \centering
% \caption{Page de consultation des données avec recherche et export}
% \end{figure}

\subsection*{E.4 Visualisation des prédictions}
\addcontentsline{toc}{subsection}{E.4 Visualisation des prédictions}

\begin{enumerate}
    \item Cliquer sur la carte "Visualisations" dans les actions rapides
    \item Sélectionner une figure dans la barre latérale gauche
    \item Consulter le graphique dans la zone principale
    \item Télécharger la figure via le lien de téléchargement
\end{enumerate}

% Espace pour figure
% \begin{figure}[h!]
% \centering
% \caption{Page des visualisations avec la figure sélectionnée}
% \end{figure}

\subsection*{E.5 Génération des prédictions}
\addcontentsline{toc}{subsection}{E.5 Génération des prédictions}

\begin{enumerate}
    \item Accéder à la page "Visualisations"
    \item Cliquer sur "Exécuter le pipeline complet"
    \item Attendre la fin de l'exécution (environ 10-15 secondes)
    \item Les nouvelles figures apparaissent automatiquement
\end{enumerate}

\subsection*{E.6 Export des résultats}
\addcontentsline{toc}{subsection}{E.6 Export des résultats}

\begin{enumerate}
    \item Accéder à la table souhaitée via le DataViewer
    \item Appliquer les filtres nécessaires
    \item Cliquer sur le format d'export désiré (CSV, Excel, JSON)
    \item Le fichier est automatiquement téléchargé
\end{enumerate}

\newpage

\section*{Annexe F : Glossaire technique}
\addcontentsline{toc}{section}{Annexe F : Glossaire technique}

\begin{longtable}{|p{3cm}|p{11cm}|}
\hline
\textbf{Terme} & \textbf{Définition} \\
\hline
API & Application Programming Interface. Interface de programmation permettant à deux applications de communiquer. \\
\hline
BI & Business Intelligence. Ensemble de technologies pour l'analyse et la visualisation de données. \\
\hline
BIAT & Banque Internationale Arabe de Tunisie. Institution bancaire tunisienne. \\
\hline
CDD & Contrat à Durée Déterminée. Type de contrat de travail temporaire. \\
\hline
CDI & Contrat à Durée Indéterminée. Type de contrat de travail permanent. \\
\hline
CRUD & Create, Read, Update, Delete. Opérations de base sur une base de données. \\
\hline
CSV & Comma-Separated Values. Format de fichier texte pour les données tabulaires. \\
\hline
DAI & Direction des Affaires Internationales. Structure de la BIAT. \\
\hline
DCO & Direction Commerciale. Structure de la BIAT. \\
\hline
DFI & Direction Financière. Structure de la BIAT. \\
\hline
DRH & Direction des Ressources Humaines. Structure de la BIAT. \\
\hline
DRS & Direction des Risques. Structure de la BIAT. \\
\hline
DSI & Direction des Systèmes d'Information. Structure de la BIAT. \\
\hline
DT & Dinar Tunisien. Monnaie officielle de la Tunisie. \\
\hline
DWH & Data Warehouse. Entrepôt de données pour l'analyse décisionnelle. \\
\hline
ETL & Extract, Transform, Load. Processus d'alimentation d'un Data Warehouse. \\
\hline
GLSI & Génie Logiciel et Systèmes d'Information. Filière d'ingénierie. \\
\hline
HTML & Hypertext Markup Language. Langage de balisage pour les pages web. \\
\hline
HTTP & Hypertext Transfer Protocol. Protocole de communication web. \\
\hline
JSON & JavaScript Object Notation. Format léger d'échange de données. \\
\hline
JWT & JSON Web Token. Standard pour l'authentification sécurisée. \\
\hline
KPI & Key Performance Indicator. Indicateur clé de performance. \\
\hline
MAE & Mean Absolute Error. Métrique d'erreur moyenne absolue. \\
\hline
MAPE & Mean Absolute Percentage Error. Métrique d'erreur relative moyenne. \\
\hline
MCD & Modèle Conceptuel de Données. Représentation abstraite des données. \\
\hline
ML & Machine Learning. Apprentissage automatique. \\
\hline
MLD & Modèle Logique de Données. Traduction du MCD en structures relationnelles. \\
\hline
MPD & Modèle Physique de Données. Implémentation technique du modèle. \\
\hline
OLAP & Online Analytical Processing. Traitement analytique en ligne. \\
\hline
OLTP & Online Transaction Processing. Traitement transactionnel en ligne. \\
\hline
PFE & Projet de Fin d'Études. Projet de fin de cycle ingénieur. \\
\hline
R² & Coefficient de détermination. Métrique de qualité de régression. \\
\hline
REST & Representational State Transfer. Style d'architecture pour les API web. \\
\hline
RMSE & Root Mean Square Error. Racine de l'erreur quadratique moyenne. \\
\hline
SGBD & Système de Gestion de Base de Données. Logiciel de gestion de données. \\
\hline
SI & Système d'Information. Ensemble des ressources informatiques. \\
\hline
SQL & Structured Query Language. Langage d'interrogation de bases de données. \\
\hline
UML & Unified Modeling Language. Langage de modélisation. \\
\hline
URL & Uniform Resource Locator. Adresse web. \\
\hline
XML & eXtensible Markup Language. Langage de balisage générique. \\
\hline
XGB & eXtreme Gradient Boosting. Modèle de Machine Learning optimisé. \\
\hline
\end{longtable}

\vspace{0.5cm}
\hrule
\vspace{0.2cm}